\originalsection{作业 7: 功与动能}
\begin{exercise}
原题 1. 质量 4.0 kg 物体从静止滑下长 3.0 m、倾角 $30^\circ$ 的斜面, 动摩擦 0.2; 到底后在动摩擦 0.3 的水平面滑到停止. (a) 斜面摩擦功及符号. (b) 重力功及符号. (c) 坡底动能. (d) 水平面摩擦功的符号表达式, 滑行距离记 $d$. (e) 求 $d$.
\end{exercise}
\begin{solution}
(a) $W_f=-\mu_1mgl\cos\theta=-20.39$ J. (b) $W_g=mgl\sin\theta=58.86$ J. (c) $K_b=38.47$ J. (d) $W_{f2}=-\mu_2mgd$. (e) 停止要求 $K_b=\mu_2mgd$, 故 $d=l(\sin\theta-\mu_1\cos\theta)/\mu_2=3.268$ m. 法向力在各段垂直位移, 不做功.
\end{solution}
\begin{exercise}
原题 2. 块 A 质量 $m_A$、速率 $v_A$ 在光滑面撞上静止块 B 质量 $m_B$ 并粘住. 组合体从 $x=0$ 经过粗糙段 $0\le x\le d$, 动摩擦系数 $\mu_k(x)=bx$, $b>0$. 到 $d$ 后进入光滑面并压缩原本未伸缩的弹簧 $k$. 求首次停止时弹簧压缩量.
\end{exercise}
\begin{solution}
碰后 $V=m_Av_A/(m_A+m_B)$. 粗糙段耗散 $\int_0^dbx(m_A+m_B)g\dd x=b(m_A+m_B)gd^2/2$. 剩余动能等于 $kX^2/2$, 所以
\begin{equation}
X=\sqrt{\frac{m_A^2v_A^2/(m_A+m_B)-b(m_A+m_B)gd^2}{k}}.
\end{equation}
根号项必须非负. 若在粗糙段已停止, 达不到弹簧, 原题的压缩情形不成立. 模型把两块视为同一集中体经过粗糙段.
\end{solution}
\begin{exercise}
原题 3. 倾角 $\phi$ 的斜面上, 质量 $m$ 物体由长 $R$ 的绳系住, 在斜面内作圆周运动. 底点 A 速率 $v_0$, 动摩擦 $\mu$. (a) 到顶点 B 的半圆过程摩擦做功多少? (b) B 处张力多少?
\end{exercise}
\begin{solution}
面法向力为 $mg\cos\phi$, 半圆路程 $\pi R$. (a) $W_f=-\mu mg\cos\phi\,\pi R$. 升高 $2R\sin\phi$, 故 $v_B^2=v_0^2-4gR\sin\phi-2\pi\mu gR\cos\phi$. (b) 顶点沿圆内方向有 $T+mg\sin\phi=mv_B^2/R$, 所以 $T=mv_0^2/R-5mg\sin\phi-2\pi\mu mg\cos\phi$. 需全过程绳绷紧; 算得顶点负张力时不能实现假设圆轨道.
\end{solution}
